\hypertarget{ux111iux1ec1u-kiux1ec7n-tiuxean-quyux1ebft}{%
\section{Điều kiện tiên
quyết}}

Trước khi bắt đầu workshop này, hãy đảm bảo bạn đã cài đặt và cấu hình
các công cụ sau.

\hypertarget{tuxe0i-khoux1ea3n-quyux1ec1n-cux1ea7n-thiux1ebft}{%
\subsection{Tài khoản \& Quyền cần
thiết}}

\hypertarget{tuxe0i-khoux1ea3n-aws}{%
\subsubsection{Tài khoản AWS}}

\begin{itemize}
\tightlist
\item
  \textbf{Tài khoản AWS} với quyền quản trị hoặc quyền cho các dịch vụ:

  \begin{itemize}
  \tightlist
  \item
    VPC, EC2, Subnets, Route Tables, Internet Gateway
  \item
    RDS (Aurora PostgreSQL Serverless v2)
  \item
    ECS (Fargate, Clusters, Task Definitions, Services)
  \item
    ECR (Repositories)
  \item
    Lambda (Functions, Layers, Permissions)
  \item
    SQS (Queues)
  \item
    EventBridge (Rules, Schedules, Targets)
  \item
    API Gateway (REST APIs)
  \item
    Bedrock (Truy cập model: \texttt{amazon.titan-embed-text-v2:0})
  \item
    IAM (Roles, Policies)
  \item
    CloudWatch (Log Groups, Metrics)
  \item
    CloudFormation (Stacks)
  \end{itemize}
\end{itemize}

\begin{quote}
\textbf{Mẹo:} Sử dụng IAM user có \texttt{AdministratorAccess} cho
workshop, hoặc đảm bảo role của bạn có các quyền trên.
\end{quote}

\hypertarget{api-keys-buxean-ngouxe0i-cho-rag}{%
\subsubsection{API Keys bên ngoài (cho
RAG)}}

\begin{itemize}
\tightlist
\item
  \textbf{Groq API Key} --- Lấy từ
  \href{https://console.groq.com/}{console.groq.com} (Có gói miễn phí)
\item
  \textbf{Google Gemini API Key} --- Lấy từ
  \href{https://aistudio.google.com/}{aistudio.google.com} (Có gói miễn
  phí)
\end{itemize}

\hypertarget{cuxf4ng-cux1ee5-bux1eaft-buux1ed9c}{%
\subsection{Công cụ bắt buộc}}

\hypertarget{aws-cli-v2}{%
\subsubsection{1. AWS CLI v2}}

\emph{(Mã nguồn chi tiết đã được lược bỏ để báo cáo ngắn gọn. Vui lòng
xem trong source code dự án)}

\textbf{Cấu hình AWS CLI:}

\begin{Shaded}
\begin{Highlighting}[]
\ExtensionTok{aws}\NormalTok{ configure}
\CommentTok{\# Nhập: AWS Access Key ID, Secret Access Key, Region (ví dụ: ap{-}southeast{-}2), Output format (json)}
\end{Highlighting}
\end{Shaded}

\hypertarget{terraform-1.5.0}{%
\subsubsection{2. Terraform \textgreater= 1.5.0}}

\emph{(Mã nguồn chi tiết đã được lược bỏ để báo cáo ngắn gọn. Vui lòng
xem trong source code dự án)}

\hypertarget{docker-docker-compose}{%
\subsubsection{3. Docker \& Docker
Compose}}

\emph{(Mã nguồn chi tiết đã được lược bỏ để báo cáo ngắn gọn. Vui lòng
xem trong source code dự án)}

\hypertarget{python-3.10}{%
\subsubsection{4. Python 3.10+}}

\emph{(Mã nguồn chi tiết đã được lược bỏ để báo cáo ngắn gọn. Vui lòng
xem trong source code dự án)}

\hypertarget{git}{%
\subsubsection{5. Git}}

\emph{(Mã nguồn chi tiết đã được lược bỏ để báo cáo ngắn gọn. Vui lòng
xem trong source code dự án)}

\hypertarget{truxecnh-soux1ea1n-thux1ea3o-muxe3-khuyux1ebfn-nghux1ecb-vs-code}{%
\subsubsection{6. Trình soạn thảo mã (Khuyến nghị: VS
Code)}}

\emph{(Mã nguồn chi tiết đã được lược bỏ để báo cáo ngắn gọn. Vui lòng
xem trong source code dự án)}

\hypertarget{thiux1ebft-lux1eadp-dux1ef1-uxe1n}{%
\subsection{Thiết lập dự án}}

\hypertarget{clone-repository}{%
\subsubsection{Clone Repository}}

\begin{Shaded}
\begin{Highlighting}[]
\FunctionTok{git}\NormalTok{ clone }\OperatorTok{\textless{}}\NormalTok{your{-}repo{-}url}\OperatorTok{\textgreater{}}\NormalTok{ AWS{-}Projects}
\BuiltInTok{cd}\NormalTok{ AWS{-}Projects}
\end{Highlighting}
\end{Shaded}

\hypertarget{muxf4i-trux1b0ux1eddng-ux1ea3o-python}{%
\subsubsection{Môi trường ảo
Python}}

\begin{Shaded}
\begin{Highlighting}[]
\ExtensionTok{python3} \AttributeTok{{-}m}\NormalTok{ venv venv}
\BuiltInTok{source}\NormalTok{ venv/bin/activate  }\CommentTok{\# Linux/macOS}
\CommentTok{\# venv\textbackslash{}Scripts\textbackslash{}activate   \# Windows}

\ExtensionTok{pip}\NormalTok{ install }\AttributeTok{{-}{-}upgrade}\NormalTok{ pip}
\ExtensionTok{pip}\NormalTok{ install }\AttributeTok{{-}r}\NormalTok{ requirements.txt}
\end{Highlighting}
\end{Shaded}

\hypertarget{biux1ebfn-muxf4i-trux1b0ux1eddng}{%
\subsubsection{Biến môi trường}}

Sao chép file mẫu và điền giá trị của bạn:

\begin{Shaded}
\begin{Highlighting}[]
\FunctionTok{cp}\NormalTok{ .env.example .env}
\end{Highlighting}
\end{Shaded}

Chỉnh sửa \texttt{.env} với thông tin xác thực:

\emph{(Mã nguồn chi tiết đã được lược bỏ để báo cáo ngắn gọn. Vui lòng
xem trong source code dự án)}

\hypertarget{biux1ebfn-terraform}{%
\subsubsection{Biến Terraform}}

Tạo \texttt{terraform.tfvars}:

\begin{Shaded}
\begin{Highlighting}[]
\FunctionTok{cat} \OperatorTok{\textgreater{}}\NormalTok{ terraform.tfvars }\OperatorTok{\textless{}\textless{} EOF}
\StringTok{db\_password     = "your\_secure\_password\_123!"}
\StringTok{db\_user         = "postgres"}
\StringTok{qdrant\_api\_key  = ""  \# Để trống cho triển khai AWS}
\StringTok{qdrant\_host     = ""  \# Để trống cho triển khai AWS}
\StringTok{model\_1\_api\_key = "gsk\_xxxxxxxxxxxxxxxxxxxx"  \# Groq}
\StringTok{model\_2\_api\_key = "gsk\_xxxxxxxxxxxxxxxxxxxx"  \# Groq backup}
\StringTok{model\_3\_api\_key = "AIza\_xxxxxxxxxxxxxxxxxxxx" \# Gemini}
\OperatorTok{EOF}
\end{Highlighting}
\end{Shaded}

\begin{quote}
\textbf{Bảo mật:} Không bao giờ commit \texttt{.env} hoặc
\texttt{terraform.tfvars} lên git. Chúng đã có trong
\texttt{.gitignore}.
\end{quote}

\hypertarget{xuxe1c-minh-thiux1ebft-lux1eadp}{%
\subsection{Xác minh thiết lập}}

Chạy script kiểm tra:

\begin{Shaded}
\begin{Highlighting}[]
\FunctionTok{make}\NormalTok{ verify}
\end{Highlighting}
\end{Shaded}

Hoặc kiểm tra thủ công từng công cụ:

\begin{Shaded}
\begin{Highlighting}[]
\ExtensionTok{aws}\NormalTok{ sts get{-}caller{-}identity  }\CommentTok{\# Nên hiển thị AWS account của bạn}
\ExtensionTok{terraform}\NormalTok{ version            }\CommentTok{\# Nên \textgreater{}= 1.5.0}
\ExtensionTok{docker}\NormalTok{ run hello{-}world       }\CommentTok{\# Nên in "Hello from Docker!"}
\ExtensionTok{python3} \AttributeTok{{-}c} \StringTok{"import scrapy; print(\textquotesingle{}Scrapy OK\textquotesingle{})"}
\end{Highlighting}
\end{Shaded}

\hypertarget{khu-vux1ef1c-aws-region}{%
\subsection{Khu vực AWS (Region)}}

Workshop này sử dụng \textbf{ap-southeast-2 (Sydney)} mặc định. Để thay
đổi: 1. Cập nhật
\texttt{provider\ "aws"\ \{\ region\ =\ "your-region"\ \}} trong
\texttt{main.tf} 2. Cập nhật \texttt{AWS\_REGION} trong \texttt{.env} 3.
Đảm bảo model Bedrock \texttt{amazon.titan-embed-text-v2:0} có sẵn trong
region của bạn (kiểm tra
\href{https://docs.aws.amazon.com/bedrock/latest/userguide/models-regions.html}{Bedrock
regions})

\hypertarget{truy-cux1eadp-model-bedrock}{%
\subsection{Truy cập Model Bedrock}}

Bật quyền truy cập model trong AWS Console: 1. Vào \textbf{Amazon
Bedrock} \(\rightarrow\) \textbf{Model access} 2. Nhấn \textbf{Manage
model access} 3. Bật \textbf{Amazon Titan Embeddings G1 - Text v2}
(\texttt{amazon.titan-embed-text-v2:0}) 4. Đợi trạng thái hiển thị
``Access granted''

\begin{center}\rule{0.5\linewidth}{0.5pt}\end{center}

\textbf{Tiếp theo:} \href{5.3-Infrastructure/}{Hạ tầng dưới dạng Code
(Terraform)}
